\subsection{有限 Hecke 范畴的范畴迹}

本小节中，我们在通常的代数几何（而非完全代数几何）中工作。我们研究 \S \ref{SS: cat of corr} 中所定义的范畴在如下情形下的表现。设 $G$ 是 $k$ 上的仿射代数群，带有自同态 $\sigma$；设 $K\subset G$ 是闭子群，且 $\sigma$ 限制在 $K$ 上仍为自同态。我们假设 $\sigma$ 诱导泛同胚，从而诱导 $G$ 的底层 \'etale 拓扑斯上的自同构。例如，若 $G$ 定义在有限域 $\bF_q$ 上，可取 $\sigma$ 为 $q$-Frobenius；另一个例子是 $\sigma=\id$。
把 $\sigma$-扭伴随作用记作
\[\Ad_\sigma: G\times G\to G,\quad \Ad_\sigma(g_1)(g_2)=\sigma(g_1)g_2 g_1^{-1}.\]
令
$$X=\frac{G}{\Ad_\sigma K},\quad Y=\frac{G}{\Ad_\sigma G},$$
分别为 $G$ 关于 $K$ 与 $G$ 的 $\sigma$-扭伴随作用的商叠。
注意 $f: X\to Y$ 紧合当且仅当 $K$ 是 $G$ 的抛物子群。

令 $C=X\times_YX$，并令 $\on{D}^{C}(X)$ 为 \S \ref{SS: cat of corr} 中由群胚 $C\rightrightarrows X$ 构造的范畴。
由引理 \ref{L: fully faithful}，它是 $G$ 上 $\Ad_\sigma G$-等变层范畴 $\on{D}(\frac{G}{\Ad_\sigma G})$ 的全子范畴，而后者是表示论中熟知的对象。
所以本小节并未构造新东西。不过，把 $G$ 上 $\Ad_\sigma G$-等变层表示成 $\on{D}^{C}(X)$ 中的对象，有助于我们重新理解表示论中许多已知的构造；而且在下一节我们转到仿射情形时，它将成为那里的构造的玩具模型——在仿射情形，$\on{D}^{C}(X)$ 的形式体系变得不可或缺。

作为热身，我们先用新的表述重写一些已知的结果与构造。

\begin{ex}
\label{Ex:springer}
设 $\sigma=\id$，$K=B$。由引理 \ref{L: fully faithful}，这是 $G$ 上 $\Ad G$-等变层范畴 $\on{D}(\frac{G}{\Ad G})$ 的全子范畴。

也可以考虑闭子叠 $\frac{B}{\Ad B}\subset \frac{G}{\Ad B}$ 及其上的常层 $\delta_{B}$。注意 $\frac{B}{\Ad B}\to \frac{G}{\Ad G}$ 沿 $G\to \frac{G}{\Ad G}$ 的基变换为
$\widetilde G:=G\times_{\Ad,B}B\to G$，通常称为 Grothendieck–Springer 变体。由经典的 Springer 理论，
$$\End_{\on{D}^{C}(X)}(\delta_{B})\cong E[W]$$
同构于 $G$ 的 Weyl 群 $W$ 的群代数。稍后我们会看到，这个 $W$ 在 $\delta_{B}$ 上的作用可由 \eqref{E: K to End} 应用于当前情形得到。更一般地，对每个 $w\in W$，可以把 $\frac{\overline{BwB}}{\Ad B}$ 上的相交上同调层看作 $\on{D}^{C}(X)$ 中的对象；在嵌入 $\on{D}^{C}(X)\subset \on{D}(\frac{G}{\Ad G})$ 下，它的幂等元将对应于 Lusztig 的幂幺特征层。

设 $U\subset B$ 为幂幺根。也可以考虑闭子叠 $\frac{U}{\Ad B}\subset \frac{G}{\Ad B}$ 及其上的常层 $\delta_{U}$。它对应于 $\on{D}(\frac{G}{\Ad G})$ 中的 Springer 层。
同样由经典的 Springer 理论，
$\End_{\on{D}^C(X)}(\delta_{U})\cong E[W]$。
\end{ex}

\begin{ex}
\label{Ex:DL}
设 $k=\bF_q$，$K=B$ 为 Borel 子群，$\sigma$ 为 $q$-Frobenius。由 Lang 定理，这是 $\on{D}^C(X)=\on{D}(\on{pt}/G(\bF_q))$ 的全子范畴。
考虑闭子叠 $\frac{B}{\Ad_\sigma B}\subset \frac{G}{\Ad_\sigma B}$。由 Lang 定理，对应
\[\frac{B}{\Ad_\sigma B}\leftarrow\frac{B}{\Ad_\sigma B}\times_{\frac{G}{\Ad_\sigma B}}\frac{B}{\Ad_\sigma B}\rightarrow \frac{B}{\Ad_\sigma B}\]
可等同于离散（Deligne–Mumford）叠之间的对应
\begin{equation}
\label{E:Hk for unit}
B(\bF_q)\backslash \on{pt}\leftarrow B(\bF_q)\backslash G(\bF_q)/B(\bF_q)\rightarrow \on{pt}/B(\bF_q).
\end{equation}
这里 $G(\bF_q)$ 与 $B(\bF_q)$ 看作 $\bF_q$ 上的离散代数群。
由此，对每个表示 $\rho: B(\bF_q)\to \GL(V)$，按通常的伴随构造有支撑在 $\frac{B}{\Ad_\sigma B}$ 上的局部系统 $\delta_\rho$，它进而定义 $\on{D}^C(X)$ 中的对象。若 $\rho$ 是平凡表示，我们把这个对象 $\delta_\rho$ 也记作 $\widetilde{\delta_e}$。
记 $H_W$ 为 $B(\bF_q)\backslash G(\bF_q)/B(\bF_q)$ 上 $E$ 值函数构成的 $E$-模，其代数结构由卷积（与 \eqref{E:conv prod alge} 比较）
\[(f*g)(x)=\sum_{y\in G(\bF_q)/B(\bF_q)} f(y)g(y^{-1}x)\]
给出。它通常称为附属于 $W$ 的、系数在 $E$ 中的 Iwahori–Hecke 代数。本质上由定义（以及例 \ref{Ex:examples of correspondences} (5)），有
\begin{equation}
\label{E:endo for unit}
\End(\widetilde{\delta_{e}})=H^{\on{op}}_W.
\end{equation}
类似地，若 $\rho=\on{Reg}$ 是 $B(\bF_q)$ 的正则表示，则
\begin{equation}
\label{E:endo for reg}
\End(\delta_{\on{Reg}})=E[G(\bF_q)]^{\on{op}},
\end{equation}
其中 $E[G(\bF_q)]$ 是 $G(\bF_q)$ 的群代数。

更一般地，对每个 $w\in W$，映射 $\overleftarrow{c}: \overrightarrow{c}^{-1}(\frac{BwB}{\Ad_\sigma B})\to \frac{G}{\Ad_\sigma B}$ 的纤维正是附属于 $w$ 的通常 Deligne–Lusztig 簇 $\on{DL}_w$。
此外有同构
\[  \frac{U}{\Ad_\sigma ((U \cap w^{-1}Uw)T^{w^{-1}\sigma})}\cong \frac{BwB}{\Ad_\sigma B},\quad u\mapsto u\dot{w},\]
其中 $T^{w^{-1}\sigma}=\{t\in T\mid \sigma(t)=wtw^{-1}\}$ 是有限环面，$\dot{w}$ 是 $w$ 到 $G$ 中 $T$ 的正规化子的提升。
因此，对该有限环面的任意特征标 $\theta$，有 $\frac{BwB}{\Ad_\sigma B}$ 上的局部系统 $\mL_{w,\theta}$；把它零延拓到 $\frac{G}{\Ad_\sigma B}$，就可看作 $\on{D}^{C}(X)$ 中的对象。
在全嵌入 $\on{D}^{C}(X)\subset \on{D}(\on{Rep}(G(\bF_q)))$ 下，它的像就是通常的 Deligne–Lusztig 表示。
\end{ex}

\begin{rmk}
\label{R:springer}
事实上可以验证，在上述两个例子中，若 $\cha E=0$，则 $\on{D}^C(X)$ 的幂等完备化就是 $\on{D}(Y)$。
\end{rmk}

下面我们说明 $\on{D}^{C}(X)$ 中态射的一般构造，其想法部分来自 \cite[\S 6]{XZ}，而后者推广了 V. Lafforgue 关于 $S$-算子的构造（\cite[\S 6]{La}）。我们将大量使用上同调对应的形式体系，\S \ref{Sec:cohomological correspondence} 中有简要回顾。

记 $X_r$ 为 $G^r$ 关于 $K^r$ 的作用
\begin{equation*}
(k_1,\ldots,k_r)\cdot (g_1,\ldots,g_r)=(\sigma(k_1)g_1k_2^{-1},k_2g_2k_3^{-1},\ldots,k_rg_rk_1^{-1}),\quad k_i\in K, \ g_i\in G
\end{equation*}
的商。特别地 $X_1=X=\frac{G}{\Ad_\sigma K}$。有时也把 $X_2$ 写作 $\frac{G\times^K G}{\Ad_\sigma K}$。
把第 $i$ 个与第 $(i+1)$ 个因子相乘的映射 $m_i:G^r\to G^{r-1}$ 诱导
\begin{equation}
\label{E: mult}
m_i: X_r\to X_{r-1},
\end{equation}
我们仍称之为乘法映射。
此外，映射 $(g_1,\ldots, g_r)\mapsto (\sigma(g_r),g_1,\ldots, g_{r-1})$ 诱导
\begin{equation}
\label{E: psigma}
\pFr: X_r\to X_r,
\end{equation}
我们称之为部分 $\sigma$-映射。注意 $(_\mathrm{p}\sigma)^r=\sigma$，这正是该名称的由来。
下面的等式是显然的。
\begin{equation*}
\label{E:mult vs pFr}
\pFr \cdot m_i=\left\{ \begin{array}{ll}m_{i-1}\cdot\pFr &  i>1\\ m_{r-1}\cdot (\pFr)^2 & i=1. \end{array}\right.
\end{equation*}

\begin{lem}
\label{L: decom of C}
存在典范同构
\[C=X\times_YX\cong X_2,\]
使得到第一个因子的左投影 $\overleftarrow{c}:X\times_YX\to X$ 等同于部分 $\sigma$-映射 \eqref{E: psigma} 与乘法映射 \eqref{E: mult} 的复合，
而到第二个因子的右投影 $\overrightarrow{c}:X\times_YX\to X$ 等同于乘法映射 \eqref{E: mult}。
\end{lem}
\begin{proof}
注意 $X\times_YX\cong K\backslash G\times^{K, \Ad_\sigma}G$，两个投影
$$\overleftarrow{c},\overrightarrow{c}:  K\backslash G\times^{K, \Ad_\sigma}G\to \frac{G}{\Ad_\sigma K}$$
分别由 $G\times G\to G, \ (g_1,g_2)\mapsto \sigma(g_1)g_2g_1^{-1}$ 与 $(g_1,g_2)\mapsto g_2$ 诱导。
而同构 $G\times G\cong G\times G,\ (g_1,g_2)\mapsto (g_2g_1^{-1}, g_1)$ 诱导
\[
 K\backslash G\times^{K, \Ad_\sigma}G \cong  X_2,\]
在此同构下，$\overleftarrow{c}$ 与 $\overrightarrow{c}$ 成为引理中所述的那些映射。
\end{proof}
